The eventual linear-independence clause is produced from the asymptotic
overlaps: diagonal overlaps tending to one and off-diagonal overlaps tending to
zero make the Gram matrix converge to the identity.

\begin{lemma}[Eventual linear independence from finite-index overlap orthonormality]
    \label{lem:bnt_finite_index_overlap_orthonormal_li}
    \lean{MPSTensor.eventually_linearIndependent_of_finite_overlap_tendsto_orthonormal}
    \leanok
    \uses{def:mpv_overlap}
    Let $I$ be a finite index set and let $(A_i)_{i \in I}$ be a finite
    family of tensors such that $O_{A_iA_i}(N) \to 1$ for each $i$, and
    $O_{A_iA_j}(N) \to 0$ whenever $i \neq j$. Then the MPV states
    $\{\ket{V^{(N)}(A_i)}\}_{i \in I}$ are linearly independent for all
    sufficiently large~$N$.
\end{lemma}

\begin{proof}\leanok
    The Gram matrix of the family $\{\ket{V^{(N)}(A_i)}\}_{i \in I}$
    converges entrywise to the identity matrix. Hence for all sufficiently
    large $N$ it is invertible, so the MPV states are linearly independent.
\end{proof}

\begin{lemma}[Eventual linear independence from overlap orthonormality]
    \label{lem:bnt_overlap_orthonormal_li}
    \lean{MPSTensor.eventually_linearIndependent_of_overlap_tendsto_orthonormal}
    \leanok
    \uses{def:mpv_overlap, lem:bnt_finite_index_overlap_orthonormal_li}
    Let $(A_j)_{j=1}^g$ be a finite family of tensors such that
    $O_{A_jA_j}(N) \to 1$ for each $j$ and $O_{A_iA_j}(N) \to 0$
    for $i \neq j$. Then the MPV states
    $\{\ket{V^{(N)}(A_j)}\}_{j=1}^g$ are linearly independent for all
    sufficiently large~$N$.
\end{lemma}

\begin{proof}\leanok
    Apply Lemma~\ref{lem:bnt_finite_index_overlap_orthonormal_li} with
    $I = \{1,\ldots,g\}$.
\end{proof}

The same conclusion is often more convenient with an explicit threshold rather
than the eventual quantifier.

\begin{theorem}[Existential BNT independence from overlap limits]%
    \label{thm:bnt_li_from_overlap_limits_exists}
    \lean{MPSTensor.exists_eventually_linearIndependent_of_overlap_tendsto_orthonormal}
    \leanok
    \uses{lem:bnt_overlap_orthonormal_li}
    If the self-overlaps of a finite block family tend to $1$ and the
    cross-overlaps of distinct blocks tend to $0$, then the MPV states are
    linearly independent for all sufficiently large system sizes.
\end{theorem}

\begin{proof}\leanok
    \uses{lem:bnt_overlap_orthonormal_li}
    This is Lemma~\ref{lem:bnt_overlap_orthonormal_li}, restated with an
    explicit threshold $N_0$.
\end{proof}

Eventual linear independence is also what converts an equality of two finite
linear combinations into equality of their coefficients; this is the mechanism
behind the coefficient identities of Chapter~\ref{ch:ft_proof}.

\begin{lemma}[Eventual coefficient extraction from linear independence]%
    \label{lem:eventual_coeff_eq_of_eventual_li}
    \lean{MPSTensor.coefficient_eventually_eq_of_eventually_linearIndependent}
    \leanok
    If a finite family of vectors is linearly independent for all sufficiently
    large $N$, and two finite linear combinations of that family are equal
    for all sufficiently large $N$, then the corresponding coefficients are
    equal for all sufficiently large $N$.
\end{lemma}

\begin{proof}\leanok
    Move the two linear combinations to one side. Linear independence forces
    each coefficient difference to vanish at every sufficiently large length.
\end{proof}

When two BNT families must be compared against each other, the same
Gram-matrix criterion applies to their disjoint union, provided the mixed
overlaps also decay. This is the form used in the exact sector matching of
Section~\ref{sec:sector_bnt_strong_match}.

\begin{lemma}[Eventual linear independence for two families]%
    \label{lem:bnt_two_family_overlap_orthonormal_li}
    \lean{MPSTensor.eventually_linearIndependent_of_two_family_overlap_tendsto_orthonormal}
    \leanok
    \uses{def:mpv_overlap, lem:bnt_finite_index_overlap_orthonormal_li}
    Let $(A_j)_{j=1}^{g_A}$ and $(B_k)_{k=1}^{g_B}$ be finite families of
    tensors. Suppose the overlaps inside each family are asymptotically
    orthonormal, and suppose every mixed overlap
    $O_{A_jB_k}(N)$ tends to~$0$. Then the combined MPV family
    \begin{align}
        \{\ket{V^{(N)}(A_j)}\}_{j=1}^{g_A}
        \cup
        \{\ket{V^{(N)}(B_k)}\}_{k=1}^{g_B}
        \notag
    \end{align}
    is linearly independent for all sufficiently large~$N$.
\end{lemma}

\begin{proof}\leanok
    Apply the finite-index Gram-matrix criterion to the disjoint union of the two
    index sets. The within-family hypotheses give the diagonal and
    off-diagonal limits inside each block, and the mixed overlap hypothesis gives
    the off-diagonal limits between the two blocks.
\end{proof}

We now isolate the separation condition that supplies the off-diagonal decay
and combine it with the canonical-form hypotheses.

\begin{definition}[No gauge-phase-equivalent distinct blocks]%
    \label{def:blocks_not_gauge_phase_equiv}
    \lean{MPSTensor.BlocksNotGaugePhaseEquiv}
    \leanok
    \uses{def:gauge_phase_equiv}
    A finite block family has no gauge-phase-equivalent distinct blocks when,
    for any two distinct indices with the same bond dimension, the corresponding
    blocks are not gauge-phase equivalent.
\end{definition}

\begin{definition}[Normal canonical form with BNT separation]\label{def:normal_cf_bnt}
    \lean{MPSTensor.IsNormalCanonicalFormBNT}\leanok
    \uses{def:is_normal_canonical_form, def:blocks_not_gauge_phase_equiv}
    A family is in \emph{normal canonical form with BNT separation} if it
    satisfies Definition~\ref{def:is_normal_canonical_form} and distinct blocks
    of the same bond dimension are not gauge-phase equivalent. The weight moduli
    are non-increasing but need not be strictly decreasing, so equal-modulus
    blocks are allowed; the grouping into a BNT is governed by gauge-phase
    equivalence, not by distinctness of moduli. Repeated equal-modulus copies of
    one basis tensor are not separate blocks here; they appear as the weights
    $\mu_{j,q}$ in $c_N^{(j)} = \sum_q \mu_{j,q}^N$ in the sector decomposition
    below.
    % Maintainer note: this is not the source two-layer BNT expansion
    %   A^i = \bigoplus_{j=1}^g \bigoplus_{q=1}^{r_j} \mu_{j,q} X_{j,q} A_j^i X_{j,q}^{-1},
    % where repeated copies stay visible through their coefficients,
    % multiplicities, and gauges. That expansion is recorded on the sector
    % decomposition (Definition~\ref{def:sector_bnt_cf}).
\end{definition}

The off-diagonal decay needed for the BNT criterion follows from the separation
condition together with the overlap-decay theorems of the preceding chapter.

\begin{theorem}[Cross-overlap decay from explicit BNT hypotheses]
    \label{thm:cross_overlap_zero_split}
    \lean{MPSTensor.cross_overlap_tendsto_zero_of_separated_bnt_data}
    \leanok
    \uses{def:mpv_overlap, def:gauge_phase_equiv}
    Let $(A_j)_{j=1}^r$ be a family in which each $A_j$ is injective and
    left-canonical, $\sum_i (A_j^i)^\dagger A_j^i = \Id$, and distinct
    equal-dimension blocks are not gauge-phase equivalent. Then
    $O_{A_jA_k}(N) \to 0$ for every pair $j \neq k$.
\end{theorem}

\begin{proof}\leanok\uses{thm:overlap_decay, thm:overlap_decay_rect}
    If the bond dimensions differ, Theorem~\ref{thm:overlap_decay_rect}
    gives the decay. If the bond dimensions agree, the non-equivalence
    hypothesis excludes gauge-phase equivalence, so
    Theorem~\ref{thm:overlap_decay} gives the decay.
\end{proof}

\begin{lemma}[BNT from explicit blockwise hypotheses]
    \label{lem:cf_bnt_is_bnt_split}
    \lean{MPSTensor.isBNT_of_separated_bnt_data}
    \leanok
    \uses{def:bnt, def:mpv_overlap}
    Let $(\mu_k, A_k)_{k=1}^r$ be a weighted block family. Assume each $A_k$ is
    injective and left-canonical, $\sum_i (A_k^i)^\dagger A_k^i = \Id$, that
    $O_{A_kA_k}(N) \to 1$ for each $k$, and that distinct equal-dimension blocks
    are not gauge-phase equivalent. Then the blocks form a basis of normal
    tensors for the block-diagonal tensor $\bigoplus_k \mu_k A_k$.
\end{lemma}

\begin{proof}\leanok
    \uses{thm:mpv_decomposition, lem:bnt_overlap_orthonormal_li, thm:cross_overlap_zero_split}
    Injective blocks are normal. The block-diagonal decomposition theorem gives
    the MPV expansion of $\bigoplus_k \mu_k A_k$ in terms of the block MPVs.
    The self-overlap limits are part of the hypotheses, and
    Theorem~\ref{thm:cross_overlap_zero_split} supplies the off-diagonal
    decay. Therefore Lemma~\ref{lem:bnt_overlap_orthonormal_li} gives the
    eventual linear independence required in Definition~\ref{def:bnt}.
\end{proof}

The normal-canonical formulation encodes normality spectrally (primitive
transfer map), the BNT definition algebraically (eventual block injectivity).

\begin{remark}[Comparing the two normality conditions]%
    \label{rem:normalcf_bnt_gap}
    The spectral normality of Definition~\ref{def:is_normal_canonical_form} and
    the algebraic normality of Definition~\ref{def:normal} agree: the Wielandt
    theory of Chapter~\ref{ch:wielandt} supplies the implication from
    primitivity to eventual block injectivity used here.
\end{remark}

\begin{lemma}[Restricted normal-CF-BNT hypotheses give a BNT]
    \label{lem:normal_cf_bnt_is_bnt}
    \lean{MPSTensor.IsNormalCanonicalFormBNT.isBNT}
    \leanok
    \uses{def:normal_cf_bnt, def:bnt}
    If a family is in normal canonical form with BNT separation and each block
    is also normal in the algebraic sense, then the blocks form a basis of
    normal tensors for the associated block-diagonal tensor. The blocks here are
    distinct representatives.
    % Scope restriction (ft_one_copy_scope_restriction): this form carries no
    % repeated-copy multiplicities; the missing multiplicity argument is
    % recorded in docs/paper-gaps/ft_one_copy_scope_restriction.tex.
\end{lemma}

\begin{proof}\leanok\uses{rem:normalcf_bnt_gap}
    The normal-canonical hypotheses already give the block-diagonal MPV
    expansion, the self-overlap normalization, and the eventual linear
    independence coming from irreducible trace-preserving overlap decay.
    Supplying algebraic normality for each block adds the remaining clause in
    Definition~\ref{def:bnt}.
\end{proof}
